\documentclass[10pt,a4paper]{amsart}
\usepackage[margin=19mm]{geometry}
\usepackage{amsmath,amssymb,mathtools}
\usepackage[scr=rsfs,cal=euler]{mathalfa}
\usepackage{xfrac}
\usepackage[unicode,hidelinks,bookmarks=false]{hyperref}
\newcommand{\ie}{i.e.\ }
\newcommand{\cf}{cf.\ }
\newcommand{\resp}{respectively\ }
\newcommand{\Subsection}{\subsection}
\providecommand{\nobreakdash}{\nobreak\hbox{-}}
\setlength{\emergencystretch}{2em}
\multlinegap0pt
\usepackage{bm}
\usepackage{bbm}
\usepackage{dsfont}
\usepackage{xspace}
\usepackage{cancel}
\usepackage{mathtools}
\usepackage{amsthm}
\makeatletter
\@ifundefined{thm}{\newtheorem{thm}{Thm}}{}
\@ifundefined{lemma}{\newtheorem{lemma}[thm]{Lemma}}{}
\@ifundefined{propo}{\newtheorem{propo}[thm]{Proposition}}{}
\makeatother
\makeatletter
\newcommand{\Gammaleft}{\ensuremath{\Gamma_\lambda^{\text{left}}}\xspace}
\newcommand{\Gammaright}{\ensuremath{\Gamma_\lambda^{\text{right}}}\xspace}
\newcommand{\T}{\ensuremath{\mathcal{T}}\xspace}
\newcommand{\est}{\text{\,e.s.t.}\xspace}
\newcommand{\hatN}{\ensuremath{\hat{\nabla}}}
\newcommand{\hatT}{\ensuremath{{\hat{T}}}}
\newcommand{\hatb}{\ensuremath{\hat{b}}}
\newcommand{\haty}{\ensuremath{{\hat{y}}}}
\@ifundefined{proof}{\newenvironment{proof}{\proofstart}{\proofend}}{}
\newcommand{\rz}{\mathbb{R}}
\newcommand{\summL}[1][{}]{\ensuremath{\sum_{\lambda\in\Lambda_{#1}}}}
\newcommand{\summLA}[1][{}]{\ensuremath{\sum_{\lambda\in\Lambda_{#1}}\alpha_\lambda}}
\makeatother
\title[Error estimates for the patch bubble method for convection-d]{Error estimates for the patch bubble method for convection-dominated channel flow problem}
\author{Eberhard B\"ansch, Pedro Morin, Itat\'i Zocola}
\date{Source: arXiv:2604.20025v1 (CC BY 4.0)}
\begin{document}
\maketitle

\noindent\emph{Source and licence.} Error estimates for the patch bubble method for convection-dominated channel flow problem. Version arXiv:2604.20025v1 (CC BY 4.0), distributed under the Creative Commons Attribution 4.0 International licence, as recorded in the arXiv licence metadata for this exact version and shown on the abstract page https://arxiv.org/abs/2604.20025v1. Full rights notice: licenses/arxiv-2604.20025-CC-BY-4.0.txt.\par\smallskip

\noindent\textit{Editorial scope.} Counted unit: the Lemma whose source label is \texttt{Lemma:12} in the article (source0.tex lines 1621--1628, source label \texttt{Lemma:12}), delivered together with the article's own complete original proof of it (source0.tex lines 1630--1759, one \texttt{proof} environment of 130 lines). No mathematics is written, repaired or re-derived: statement and proof are the article's own text line for line. Required context retained uncounted: eq:defBubbCorr (lines 891--899); Lemma:boundBubbleCorr (lines 1420--1427); Temam:zetaBubble (lines 2327--2338); propo: bh\_inf\_bound (lines 3116--3125); Propo:AngleSpace2 (lines 3208--3216). Every definition, display and label of those lines is retained unchanged. Omitted from this package and not redistributed: 1--890, 900--1419, 1428--1620, 1629--1629, 1760--2326, 2339--3115, 3126--3207, 3217--3235. Nothing inside a retained excerpt is omitted or altered: every retained body is delivered exactly as the article's lines are written, so no omission receipt of any kind accompanies this package. Citations: the retained excerpts carry no citation command at all, so no bibliography entry of the article is transcribed and no reference is redistributed. Reference adaptation: none was needed and none was made; every reference inside the delivered unit resolves inside it, so no unresolved reference or citation is produced. Printed slips kept literal: no printed word, symbol, hypothesis or number of the counted statement or of its proof was altered. Layout and class: the article's own class is \texttt{article} with packages including fontenc, inputenc, multirow, amsfonts, mathtools, amssymb, graphicx, wrapfig, tikz, animate, subcaption, lmodern, fancyhdr, amsmath; the wrapper is the lane's pinned \texttt{amsart} editorial wrapper at 10pt on A4, which restates the theorem-like environments the retained text needs and adds the article's own macro definitions that the retained text needs (\texttt{\textbackslash Gammaleft}, \texttt{\textbackslash Gammaright}, \texttt{\textbackslash T}, \texttt{\textbackslash est}, \texttt{\textbackslash hatN}, \texttt{\textbackslash hatT}, \texttt{\textbackslash hatb}, \texttt{\textbackslash haty}, \texttt{\textbackslash rz}, \texttt{\textbackslash summL}, \texttt{\textbackslash summLA}), with extra wrapper packages (bm, bbm, dsfont, xspace, cancel, mathtools). The delivered numbering is the unit's own. Assets: no retained span contains a figure environment, a graphics-inclusion command, a PDF-page inclusion, a table or a TikZ picture; no image file is copied into this package, the package declares no asset dependency and no image file is redistributed with it. Licence: version arXiv:2604.20025v1 of this article is distributed under the Creative Commons Attribution 4.0 International licence, as recorded in the arXiv licence metadata for this exact version and shown on the abstract page https://arxiv.org/abs/2604.20025v1. A full rights notice is delivered at licenses/arxiv-2604.20025-CC-BY-4.0.txt. Characters: every retained excerpt is pure ASCII, so the delivered engine, XeLaTeX with its default Latin Modern text font, reports no missing character.
\begin{equation}\label{eq:defBubbCorr}
    \theta_B := \summLA \theta_\lambda,
    \quad
    \varphi_B := \summLA \varphi_\lambda,
    \quad
    \xi_B := \summLA \xi_\lambda,
    \quad
    \zeta_B := \summLA \zeta_\lambda,
\end{equation}

\begin{lemma}
\label{Lemma:boundBubbleCorr}
Let $\varphi_B, \xi_B, \zeta_B$ denote the correctors corresponding to $u_B$ given by \eqref{eq:defBubbCorr}.
Then,
$$
\|\varphi_B\|_\epsilon  +  \|\xi_B\|_\epsilon + \|\zeta_B\|_\epsilon  \lesssim  (\epsilon/h)^{1/4}.
$$
\end{lemma}

\begin{lemma}\label{Temam:zetaBubble}
Let $\lambda \in \Lambda$ and let $\zeta
_\lambda$ be an ordinary corner corrector of the bubble $\psi_\lambda$. Then
$$
\|\hat{\zeta}_\lambda e^x\|_{\hat{\Omega}_\lambda} \lesssim h \,\hat{\epsilon}^{3/4},
\quad \text{ and }\quad
\|\hat{y}\hat{\zeta}_\lambda e^x\|_{\Gammaleft} \lesssim h \,\hat{\epsilon}^{3/4},
$$
where \(\Gammaleft := \left\{ (\hat{x}, \hat{y}) \in \partial \hat{\Omega}_\lambda \,:\, \hat{x} = \hat x_{\min} \right\},
\)
with $\hat x_{\min}$ denoting the minimal $\hat{x}$-coordinate of $\hat{\Omega}_\lambda$.
\end{lemma}

\begin{propo}\label{propo: bh_inf_bound}
  Let $\hatb_h= \hatb_1 + \alpha \hatb_2 + \beta \hatb_T$ on $\hatT$ for some $\alpha,\beta\in \rz$. Then
  \begin{equation*}
  \frac{h^2(1+|\alpha|+|\beta|)^2}{\sqrt{\hat{\epsilon}}}
  \lesssim \|\hatN \hatb_h \|_\hatT^2 
  \quad\text{and}\quad
      \left\| \frac{\hat{b}_1(0,\haty)}{\haty} \right\|_{\{\hat{x}=0\}}^2 \lesssim
      \|\hatN \hatb_h \|_\hatT^2.
  \end{equation*}
\end{propo}

\begin{propo}
\label{Propo:AngleSpace2}
There is $\hat{\epsilon}_0, 0 < \hat{\epsilon}_0 < 1$ such that for
${v_h = q_h + b_h}\in V_h$ with $q_h \in V_L$  and $b_h \in V_B$ it holds
\begin{equation*}
    \|\nabla q_h\|_T  + \|\nabla b_h\|_T \lesssim \|\nabla v_h\|_T ,
\end{equation*}
for all $T\in \T$ and $0<\epsilon / h <hat{\epsilon}_0$.
\end{propo}

\begin{lemma}\label{Lemma:12}
Let $\zeta_B$ denote the $\zeta$-corrector corresponding to $u_B$ given by \eqref{eq:defBubbCorr}. There is $\hat{\epsilon}_0, 0 < \hat{\epsilon}_0 < 1$
such that
$$
 a(\zeta_B,v_h) \lesssim (\epsilon/h)^{1/4} \|v_h\|_\epsilon, \qquad \text{for all $v_h \in V_h$}
$$
for all $0<\epsilon/h <\hat{\epsilon}_0$.
\end{lemma}

\begin{proof}
Recall that
\[\zeta_B = \summLA[B]\zeta_\lambda+ \summLA[T]\zeta_\lambda+
\summLA[L]\zeta_\lambda,\]
and let us focus first on the integrals corresponding to the indices in $\Lambda_B$. 
Let $v_h\in V_h$, then, analogously as before,
\begin{equation}\label{Eq:79}
a\big(\summLA[B]\zeta_\lambda,v_h\big)  \lesssim 
\big\|\zeta_B\big\|_\epsilon \big\|v_h\big\|_\epsilon + \summLA[B] \Big|\int_{\Omega_\lambda} \partial_x \zeta_\lambda v_h e^x\Big|.
\end{equation}
Like above, by scaling and integration by parts we get
\begin{equation}\label{eq:xi_boundary}
\begin{aligned}
   \int_{\Omega_\lambda} \partial_x \zeta_\lambda v_h e^x 
   %& =  h\int_{\hat{\Omega}_\lambda} \partial_{\hat{x}} \hat{\zeta}_\lambda \hat{v}_h e^x \\
  &= -h\int_{\hat{\Omega}_\lambda} \hat{\zeta}_\lambda \partial_{\hat{x}}(\hat{v}_h e^x)
  - h\int_{\Gammaleft} \hat{\zeta}_\lambda \hat{v}_h e^x  + h\int_{\Gammaright} \text{\est } \hat{v}_h,
\end{aligned}
\end{equation}
where we consider $\Gammaleft$ and $\Gammaright$, the left and right boundary of the domain $\hat{\Omega}_\lambda$. That is,
\begin{equation*}
 \Gammaleft := \left\{ (\hat{x}, \hat{y}) \in \partial \hat{\Omega}_\lambda \,:\, \hat{x} = x_{\min} \right\}\text{\quad and \quad}
 \Gammaright := \left\{ (\hat{x}, \hat{y}) \in \partial \hat{\Omega}_\lambda \,:\, \hat{x} = x_{\max} \right\},
\end{equation*}
where $x_{\min}$ and $x_{\max}$ denotes the minimal and maximal $\hat{x}$-coordinate of $\hat{\Omega}_\lambda$, correspondingly.

Since $\partial_{\hat{x}} e^x = h e^x$, the first term of the last line of~\eqref{eq:xi_boundary} can be estimated as
\begin{equation}\label{Eq:77}
\begin{aligned}
\Big|h\int_{\hat{\Omega}_\lambda} \hat{\zeta}_\lambda \partial_{\hat{x}}(\hat{v}_h e^x)\Big|
 &\leq \Big|h\int_{{\hat{\Omega}_\lambda}} \hat{\zeta}_\lambda \partial_{\hat{x}}(\hat{v}_h) e^x\Big|
+ \Big|h^2\int_{\hat{\Omega}_\lambda} \hat{\zeta}_\lambda \hat{v}_h e^x\Big|
\\
&\le \|\hat{\zeta}_\lambda\|_{\hat{\Omega}_\lambda}
\big( h\big\|\partial_{\hat{x}} \hat{v}_h\big\|_{\hat{\Omega}_\lambda} +h^2 \big\| \hat{v}_h\big\|_{\hat{\Omega}_\lambda}\big)
\\
&\lesssim h \, \hat{\epsilon}^{3/4} \, \big(h\big\|\partial_{\hat{x}} \hat{v}_h\big\|_{\hat{\Omega}_\lambda} +h^2 \big\| \hat{v}_h\big\|_{\hat{\Omega}_\lambda}\big)
\\
&= h^2\hat{\epsilon}^{3/4} \big(\big\|\partial_x v_h\big\|_{\Omega_\lambda} +\big\|v_h\big\|_{\Omega_\lambda}\big),
\end{aligned}
\end{equation}
where in the last inequality, we have used Lemma~\ref{Temam:zetaBubble}.

The integral over \Gammaleft in~\eqref{eq:xi_boundary}   can be handled as follows:
Let ${v_h = q_h + b_h}$ with $q_h \in V_L$ being piecewise bilinear and $b_h \in V_B$  a bubble function.

Using Lemma~\ref{Temam:zetaBubble}, we obtain:
\begin{equation*}
\begin{aligned}
\Big|h\int_{\Gammaleft} \hat{\zeta}_\lambda \hat{v}_h e^x\Big| &= 
    \Big| h\int_{\Gammaleft} \hat{y} \hat{\zeta}_\lambda e^x \frac{\hat{v}_h}{\hat{y}} \Big|
  \lesssim h^2\hat{\epsilon}^{3/4} \bigg\| \frac{\hat{v}_h}{\hat{y}}\bigg\|_{\Gammaleft} 
  \\
  & \le h^2\hat{\epsilon}^{3/4} \Big(\bigg\| \frac{\hat{q}_h}{\hat{y}}\bigg\|_{\Gammaleft}+ \bigg\| \frac{\hat{b}_h}{\hat{y}}\bigg\|_{\Gammaleft}\Big).
\end{aligned}
\end{equation*}

For the first term, using Hardy's one-dimensional inequality and the equivalence of norms on finite dimensional spaces% 
%in the second to last inequality
, we obtain that:
\begin{equation*}
 \begin{aligned}
  h^2\hat{\epsilon}^{3/4} \left\| \frac{\hat{q}_h}{\hat{y}}\right\|_{\Gammaleft}
 &\lesssim  h^2\hat{\epsilon}^{3/4} \| \partial_{\hat{y}} \hat{q}_h\|_{\Gammaleft} 
 \lesssim h^2\hat{\epsilon}^{3/4} \| \hat{\nabla} \hat{q}_h\|_{\hat{\Omega}_\lambda}
 % \\
 % & 
 \lesssim  h^2\hat{\epsilon}^{3/4} \| \nabla q_h\|_{\Omega_\lambda}.
 \end{aligned}   
\end{equation*}

For the second term, due to the support of the bubbles that define $V_B$, we have $\hat{b}_h = c \, \hat{b}_{S_0}$ at $\hat{x} = x_{\min}$, where $\hat{b}_{S_0}$ denotes the patch bubble associated with the side $\hat{x} = x_{\min}$, for some constant $c$. Applying Proposition~\ref{propo: bh_inf_bound}, we obtain:

\begin{equation*}
      h^2\hat{\epsilon}^{3/4} \bigg\| \frac{\hat{b}_h}{\hat{y}}\bigg\|_{\Gammaleft}
      \lesssim 
      h^2\hat{\epsilon}^{3/4} \| \nabla  \hat{b}_h\|_{\hat{\Omega}_\lambda}.
\end{equation*}

Then, using Proposition \ref{Propo:AngleSpace2} 
one finds $\hat{\epsilon}_0, 0 < \hat{\epsilon}_0 < 1$ such that
\begin{equation}\label{Eq:78}
    |h\int_{\Gammaleft} \hat{\zeta}_\lambda \hat{v}_h e^x| \lesssim
     h^2\hat{\epsilon}^{3/4} \| \nabla v_h\|_{\Omega_\lambda}.
\end{equation}
for all $0<\epsilon/h < \hat{\epsilon}_0$.
For the last term in~\eqref{eq:xi_boundary}, by equivalence of norms, we obtain
\begin{equation}
\label{Eq:80}
    h\int_{\Gammaright} \text{\est } \hat{v}_h
    \leq
    h\est \int_{\hat{\Omega}_{\lambda}}\hat{v}_h,
\end{equation}

By combining~\eqref{Eq:79}, Lemma~\ref{Lemma:boundBubbleCorr}, and equations~\eqref{Eq:77},\eqref{Eq:78},\eqref{Eq:80}, using that $|\alpha_\lambda |\le 2N$, $\#\Lambda_B \lesssim N$, and $N = 1/h$, we conclude that

$$
\begin{aligned}
a\Big(\summLA[B] \zeta_\lambda,v_h\Big) 
    &\lesssim 
    \hat{\epsilon}^{1/4} \|v_h\|_\epsilon
 +  \summLA[B] h^2 \hat{\epsilon}^{3/4} (\|\nabla v_h\|_{\Omega_\lambda} +  \|v_h\|_{\Omega_\lambda}) \\
 %   &= \hat{\epsilon}^{1/4} \|v_h\|_\epsilon
 % +  \summL[B] \sqrt{h}  \hat{\epsilon}^{1/4} \sqrt{\epsilon} (\|\nabla v_h\|_{T_{i,1}}  + \|v_h\|_{T_{i,1}}) \\
  & \lesssim
  \hat{\epsilon}^{1/4} \|v_h\|_\epsilon
 +  \hat{\epsilon}^{1/4} \Big(\summL[B] h\Big)^{1/2}  \Big(\summL[B] \epsilon \|\nabla v_h\|_{\Omega_\lambda}^2 + \|v_h\|_{\Omega_\lambda}^2\Big)^{1/2}  \\
   & \lesssim \hat{\epsilon}^{1/4} \|v_h\|_\epsilon.
\end{aligned}
$$

Analogously we can prove the same bound for $a\Big(\summLA[B] \zeta_\lambda,v_h\Big) $, whence
$$
\begin{aligned}
a\Big(\summLA[B] \zeta_\lambda+\summLA[T] \zeta_\lambda,v_h\Big) 
   & \lesssim \hat{\epsilon}^{1/4} \|v_h\|_\epsilon.
\end{aligned}
$$

We finally consider the outflow boundary. That is, we now consider the indices $\lambda\in\Lambda_L$, and
proceed again as for the parabolic boundary layer. The estimate is even simpler, 
since the boundary term
$$
h\int_{\Gammaleft} \hat{\zeta}_\lambda \hat{v}_h e^x = 0,
$$ because $v_{h|\partial\Omega}=0$.
%
% It remains to bound the correctors in the two outflow boundary corners. 
% Observing that $u_B$ consists of bubble functions like in the parabolic and outflow boundary layer, the
% analysis is the same as above.
\end{proof}

\end{document}
